\documentclass[11pt]{amsart}

\usepackage[T1]{fontenc}
\usepackage{lmodern}
\usepackage{microtype}
\usepackage{mathtools,amssymb,amsthm}
\usepackage[hidelinks]{hyperref}

\hypersetup{
  pdftitle={A Length-14 Negative Answer to Alderson's q=4 Optimality Question}
}

\newtheorem{theorem}{Theorem}
\newtheorem{corollary}[theorem]{Corollary}
\theoremstyle{remark}
\newtheorem{remark}[theorem]{Remark}

\newcommand{\F}{\mathbb F}
\newcommand{\PG}{\mathrm{PG}}
\newcommand{\wt}{\operatorname{wt}}

\title[A length-14 negative answer]
{A Length-14 Negative Answer to Alderson's
$q=4$ Optimality Question}
\date{}

\subjclass[2020]{94B05, 51E20}
\keywords{additive code, code extension, maximal code, partial spread,
Baer subgeometry}

\begin{document}

\begin{abstract}
Alderson constructed an extendable, additively maximal additive
$(112,2,104)_{16/4}$-code and asked whether length $112$ is optimal for
$q=4$. A classical spread of
$\PG(3,4)\setminus\PG(3,2)$ yields a faithful additive MDS
$(14,2,13)_{16/4}$-code that is extendable as a $16$-ary code but admits no
$\F_4$-additive extension. Hence the answer is negative and the minimum
length in the $q=4$ clause of Alderson's Problem~9.2 is at most $14$.
Problem~9.2 itself notes that any sub-multiset of its $112$ external lines
whose maximal-fold points still meet every line yields a shorter example;
that condition governs additive maximality alone, and what is supplied here
is a sub-multiset for which extendability, the other requirement, also
holds. We do not determine the global minimum.
\end{abstract}

\maketitle

\section{Result}

Alderson's Theorem~6.6 gives, for $t=2$, an extendable, additively maximal
additive $(112,2,104)_{16/4}$-code, whose $112$ coordinates are indexed by
the lines of $\Pi=\PG(3,4)$ external to a scattered $\F_2$-linear set.
Problem~9.2 of \cite{Alderson} then reads, in full:
\begin{quote}
Extend Theorem~6.6 to non-square, non-prime $q$. Further, determine the
minimum length of an extendable, additively maximal additive
$(n,2,d)_{q^2/q}$-code for square $q$: is $n=112$ optimal for $q=4$? (Any
sub-multiset of the external lines whose set of maximal-fold points still
meets every line of $\Pi$ yields a shorter example.)
\end{quote}
Theorem~\ref{thm:main} below answers the final clause negatively, at length
$14$; the first clause and the minimum for general square $q$ are untouched.
Remark~\ref{rem:relation} records what the closing parenthetical already
supplies and what it does not.

The geometric input is classical: van Dam \cite{vanDam} classified the
spreads of $\PG(3,4)\setminus\PG(3,2)$ into three orbits under the
collineation group $\mathrm{P\Gamma L}(4,4)$, and Mellinger \cite{Mellinger}
later obtained the corresponding three mixed partitions consisting of one Baer
subgeometry and fourteen lines.

\begin{theorem}\label{thm:main}
There exists a faithful additive MDS $(14,2,13)_{16/4}$-code $C$ that is
extendable as a $16$-ary code but is additively maximal. Its Hamming weight
enumerator is
\[
  W_C(z)=1+210z^{13}+45z^{14}.
\]
\end{theorem}

\begin{proof}
Let $\F_2\subset\F_4\subset\F_{16}$, choose
$\omega\in\F_4$ with $\omega^2+\omega+1=0$, put $E=\F_4^4$, and let
$U=\F_2^4\le E$. In $\Pi=\PG(E)=\PG(3,4)$, define the standard Baer
subgeometry
\[
  B=\{\langle u\rangle_{\F_4}:0\ne u\in U\}\cong\PG(3,2).
\]
Thus $|B|=15$. The $\F_4$-line through two distinct points of $B$ meets
$B$ in the three points of the corresponding $\F_2$-subline. In
particular, $B$ contains no $\F_4$-line.

We first give an explicit representative of the classical complement
spreads classified in \cite{vanDam,Mellinger}. Let
\[
 T=\begin{pmatrix}
     0&0&1&0\\
     1&0&1&0\\
     0&1&0&0\\
     0&0&0&1
   \end{pmatrix}\in\mathrm{GL}(4,2)\le\mathrm{GL}(4,4).
\]
The upper-left block has minimal polynomial $X^3+X+1$, so $T$ has order
$7$. Put
\[
\begin{aligned}
 u&=(1,0,0,\omega)^{\mathsf T},
 &v&=(0,1,\omega^2,0)^{\mathsf T},
 &L&=\langle u,v\rangle_{\F_4},\\
 u'&=(1,0,0,\omega^2)^{\mathsf T},
 &v'&=(0,1,\omega,0)^{\mathsf T},
 &L'&=\langle u',v'\rangle_{\F_4}.
\end{aligned}
\]
The five points of $\PG(L)$ are represented by
$u,v,u+v,u+\omega v,u+\omega^2v$. These representatives are normalized
at their first nonzero coordinate, and each has a coordinate outside
$\F_2$; hence none lies in $B$. The same holds for $\PG(L')$, its
Frobenius conjugate. Since $T$ has entries in $\F_2$, every line
$\PG(T^jL)$ and $\PG(T^jL')$ is disjoint from $B$.

For vectors $x_1,\ldots,x_4$, write
$\det[x_1,x_2,x_3,x_4]$ for the determinant of the matrix with those
columns. A direct calculation gives
\begin{align*}
 \bigl(\det[u,v,T^ku,T^kv]\bigr)_{k=1}^{6}
   &=(\omega,1,\omega^2,\omega^2,1,\omega),\\
 \bigl(\det[u',v',T^ku',T^kv']\bigr)_{k=1}^{6}
   &=(\omega^2,1,\omega,\omega,1,\omega^2),\\
 \bigl(\det[u,v,T^ku',T^kv']\bigr)_{k=0}^{6}
   &=(1,\omega,\omega,\omega^2,\omega,\omega^2,\omega^2).
\end{align*}
All entries are nonzero. Since intersections depend only on relative
powers of $T$, these three calculations show that the fourteen lines
\[
  \ell_{j+1}=\PG(T^jL),\qquad
  \ell_{j+8}=\PG(T^jL')\qquad(0\le j\le6)
\]
are pairwise disjoint. They contain $14\cdot5=70=85-15$ points, all
outside $B$, and therefore
\begin{equation}\label{eq:partition}
  \Pi=B\mathbin{\dot\cup}\ell_1\mathbin{\dot\cup}\cdots
       \mathbin{\dot\cup}\ell_{14}.
\end{equation}

Write $\ell_i=\PG(V_i)$, where $V_i\le E$ has $\F_4$-dimension $2$.
Identifying $\F_{16}$ with a two-dimensional $\F_4$-space, choose an
$\F_4$-linear epimorphism
\[
  \rho_i:E\longrightarrow\F_{16},
  \qquad \ker\rho_i=V_i,
\]
and define
\[
  c(a)=\bigl(\rho_1(a),\ldots,\rho_{14}(a)\bigr),
  \qquad C=\{c(a):a\in E\}.
\]
Since the lines are disjoint,
\[
  \ker c=\bigcap_{i=1}^{14}V_i
  \subseteq V_1\cap V_2=\{0\}.
\]
Hence $c$ is injective and $|C|=4^4=16^2$. Each coordinate map is
surjective, so $C$ is faithful.

For $a\ne0$, the zero coordinates of $c(a)$ are exactly the indices $i$
for which $\langle a\rangle_{\F_4}\in\ell_i$. By
\eqref{eq:partition}, $c(a)$ has weight $14$ when
$\langle a\rangle_{\F_4}\in B$ and weight $13$ otherwise. Since
$|\Pi|=85$, there are $3|B|=45$ nonzero words of weight $14$ and
$3(85-15)=210$ of weight $13$. This proves the stated weight enumerator
and minimum distance. Since $13=14-2+1$, the code is MDS.

Viewed over $\F_2$, both $E/U$ and $\F_{16}$ have dimension $4$. Choose
an $\F_2$-linear isomorphism $\tau:E/U\to\F_{16}$ and set, following the
coset argument in the proof of \cite[Theorem~6.6]{Alderson},
\[
  C^+=\{\bigl(c(a),\tau(a+U)\bigr):a\in E\}
  \subseteq\F_{16}^{15}.
\]
For distinct $a,b\in E$, put $x=a-b$. If $x\in U$, then the new symbols
agree and $\langle x\rangle_{\F_4}\in B$, so
$d_H(c(a),c(b))=14$. If $x\notin U$, then the new symbols differ and
$d_H(c(a),c(b))\ge13$. Thus $d(C^+)=14$, with equality in the first
case, and puncturing the last coordinate gives $C$. Hence $C$ is
extendable; in fact, this extension is $\F_2$-linear, the subfield
phenomenon of \cite[Remark~6.7]{Alderson}. So $C$, regarded not as an
$\F_4$-linear but as an $\F_2$-linear subspace of $\F_{16}^{14}$ --- that is,
as an additive $(14,2,13)_{16/2}$-code in the general
$(n,k,d)_{q^m/q}$ notation of \cite{Alderson}, taken there with $q=2$ and
$m=4$, so that the $\F_2$-dimension is $km=8$, as it is here since
$|C|=2^8$ --- does admit an additive extension: the additive
maximality proved next is relative to the declared field of linearity
$\F_4$.

Suppose that $C$ admitted an $\F_4$-additive extension
$\widetilde C\le\F_{16}^{15}$ of minimum distance $14$. Puncturing the
new coordinate is an $\F_4$-linear bijection $\widetilde C\to C$.
Therefore, for some $\F_4$-linear map $\rho:E\to\F_{16}$,
\[
  \widetilde C=\{\bigl(c(a),\rho(a)\bigr):a\in E\}.
\]
If $\langle a\rangle_{\F_4}\notin B$, then $\wt(c(a))=13$, so
$\rho(a)\ne0$. Consequently every nonzero vector in $\ker\rho$ has
projective direction in $B$. But $\dim_{\F_4}\F_{16}=2$, and hence
\[
  \dim_{\F_4}\ker\rho\ge 4-2=2.
\]
The projectivization of a two-dimensional $\F_4$-subspace of
$\ker\rho$ is an $\F_4$-line contained in $B$, a contradiction. Thus
$C$ has no $\F_4$-additive extension.
\end{proof}

\begin{corollary}
The minimum length asked for in the $q=4$ clause of Alderson's
Problem~9.2 is at most $14$. In particular, $112$ is not optimal.
\end{corollary}

\begin{proof}
Theorem~\ref{thm:main} has exactly the required extension and additive
maximality properties.
\end{proof}

\begin{remark}\label{rem:relation}
Problem~9.2 already states that any sub-multiset of the $112$ external
lines whose set of maximal-fold points still meets every line of $\Pi$
yields a shorter example, and the code of Theorem~\ref{thm:main} is such a
sub-multiset: each of $\ell_1,\ldots,\ell_{14}$ is taken once, so by
\eqref{eq:partition} the maximal fold is $1$ and the maximal-fold set is
$\Pi\setminus B$, which meets every line of $\Pi$ because $B$ contains no
$\F_4$-line. For $t=2$ Alderson's scattered $\F_2$-linear set is the
projectivization of $\{(x,y,x^2,y^2):x,y\in\F_4\}$, an $\F_2$-subspace of
$E$ whose $\F_2$-basis is also an $\F_4$-basis of $E$, hence the image of $U=\F_2^4$
under an element of $\mathrm{GL}(4,4)$; it is therefore a Baer subgeometry
$\PG(3,2)$ of $\Pi$ (cf.\ \cite[Lemma~6.4]{Alderson}), projectively
equivalent to the set $B$ above. The lines of $\Pi$ external to a Baer
subgeometry number exactly
$112$: of the $357$ lines of $\Pi$, $112$ miss $B$, $210$ meet it in one
point and $35$ in three. That count is Alderson's $n_2=112$, so
$\ell_1,\ldots,\ell_{14}$ are fourteen of his own $112$ external lines, each
taken once; the maximal-fold set $\Pi\setminus B$ is his; and $C^+$ is his
coset construction.

What that parenthetical supplies is one of the two requirements. Its
condition, that the maximal-fold points meet every line of $\Pi$, is the
criterion for additive maximality of \cite[Corollary~3.6]{Alderson}, in which
the flats to be met are the $(km-m-1)$-flats of $\Pi$ and hence, for
$k=m=2$, its lines; it is silent on extendability, which Problem~9.2 also
demands, so a sub-multiset meeting the condition need not answer the problem. The content of
Theorem~\ref{thm:main} is accordingly the choice of sub-multiset: a
line--Baer partition of $\Pi$ makes the fourteen chosen lines pairwise
disjoint, so the code is MDS, and leaves the uncovered points equal to the
Baer subgeometry $B$, and it is that hole set which makes the $16$ cosets of
$U$ available as the fibres of a $16$-ary extension of minimum distance
$14$.
\end{remark}

\begin{remark}
The corollary answers only the final yes-or-no clause of Alderson's
Problem~9.2, refuting the optimality of $112$ for $q=4$. It does not
determine the global minimum or address the nonsquare, nonprime case, and no
lower bound is proved here. The spread in \eqref{eq:partition} is classical;
the contribution is its coding-theoretic consequence in Alderson's extension
framework. In the projective-systems dictionary used in \cite{Alderson}, a
faithful additive MDS $(n,2,n-1)_{q^2/q}$-code is a partial spread of $n$
lines of $\PG(3,q)$, and additive maximality corresponds, by
\cite[Corollary~3.6]{Alderson} again, to that partial spread being maximal; the object here is thus a
maximal partial spread of $\PG(3,4)$ of size $14$ whose holes form a Baer
subgeometry. Additive MDS codes over $\F_{16}$ that are $\F_4$-linear are
studied in \cite{BGL}, where the longest ones are classified; additive
extendability in the sense of \cite{Alderson} is not considered there.
\end{remark}

\section*{Exact verification}

Everything needed to redo the computation is printed above: the relation
$\omega^2+\omega+1=0$, the matrix $T$, the vectors $u,v,u',v'$, and the
partition \eqref{eq:partition}. From these one forms the fourteen lines
$\ell_{j+1}=\PG(T^jL)$ and $\ell_{j+8}=\PG(T^jL')$, checks the nineteen
determinants displayed in the proof, tests all $91$ pairs of lines for
disjointness, builds $c$, $C$ and $C^+$ as in the proof, recovers
$W_C(z)=1+210z^{13}+45z^{14}$ and the minimum distances $d(C)=13$ and
$d(C^+)=14$, and finally runs over the $\F_4$-linear maps
$\rho:E\to\F_{16}$ to confirm that none of them extends $C$ to minimum
distance $14$. A referee can therefore reproduce every number printed here
from this paper alone.

The exhibited object and every count printed above were recomputed on that
basis independently, on a separate machine, by a program written from the
definitions of this note; all $71$ of its checks agree with the paper. The
enumerations are exhaustive rather than sampled: all $85$ points and all
$357$ lines of $\Pi$, all $91$ pairs of the fourteen lines, all $256$ words of $C$ and of $C^+$ with all $32{,}640$
unordered pairs for each minimum distance, all $357$ two-dimensional
$\F_4$-subspaces of $E$, and, in the maximality step, all $4^8=65{,}536$
$\F_4$-linear maps $E\to\F_{16}$, the completeness of that last list being
itself checked. The run repeats the construction for other choices of
$\rho_1,\ldots,\rho_{14}$ and of $\tau$, reports failures on seven
deliberately corrupted variants of the object, and passes a positive control
in which the hole set does contain a line: replacing $\ell_1,\ldots,\ell_{14}$
by fourteen of the seventeen lines of a spread of $\PG(3,4)$ leaves as holes
the three remaining, pairwise disjoint lines, and the same exhaustive search
then accepts exactly $3\,|\mathrm{GL}(2,4)|=540$ of the $65{,}536$ maps
$\rho$, namely, for each hole line $\PG(V)$, the $180$ maps with kernel $V$.
So the negative answer for the Baer hole set is not an artifact of the
search. The program and its transcript accompany this note; nothing above
depends on them, since every quantity quoted is either
printed here or recomputable from the data just listed.

What the re-run does not do. It does not re-derive anything from
\cite{Alderson}: Theorem~6.6, the length $112$, and the definitions of
extendable and of additively maximal are taken from that paper; and the
identification in Remark~\ref{rem:relation} of Alderson's scattered
$\F_2$-linear set with a Baer subgeometry projectively equivalent to $B$ is
argued there rather than checked by the program, which counts the $112$ lines
external to $B$ itself. It does not
verify the classifications of \cite{vanDam,Mellinger}, for which the
explicit spread above is a self-contained substitute. And faithfulness is
tested only in the sense used here, that every coordinate projection is onto
$\F_{16}$. No search for a lower bound or for the global minimum was
attempted. The re-run and this paper agree; that is agreement, not a proof.

\begin{thebibliography}{9}

\bibitem{Alderson}
T.~L. Alderson,
\emph{On the maximality of additive codes},
arXiv:2607.22297v1 (2026),
\href{https://arxiv.org/abs/2607.22297}{arXiv:2607.22297}.

\bibitem{BGL}
S.~Ball, G.~Gamboa, M.~Lavrauw,
\emph{On additive MDS codes over small fields},
arXiv:2012.06183 (2020),
\href{https://arxiv.org/abs/2012.06183}{arXiv:2012.06183}.

\bibitem{vanDam}
E.~R. van Dam,
\emph{Classification of spreads of
$\PG(3,4)\setminus\PG(3,2)$},
Des. Codes Cryptogr. \textbf{3} (1993), no.~3, 193--198,
\href{https://doi.org/10.1007/BF01388480}
{doi:10.1007/BF01388480}.

\bibitem{Mellinger}
K.~E. Mellinger,
\emph{A note on line--Baer subspace partitions of $\PG(3,4)$},
J. Geom. \textbf{72} (2001), no.~1--2, 128--131,
\href{https://doi.org/10.1007/s00022-001-8574-0}
{doi:10.1007/s00022-001-8574-0}.

\end{thebibliography}

\end{document}